\section{Casi d'uso implementati}
In questa seconda fase di sviluppo, l'attenzione si è concentrata sull'estensione delle funzionalità di gestione dei viaggi, con particolare riguardo alla condivisione e alla modifica granulare degli itinerari. Sono stati implementati il catalogo pubblico (con la relativa funzione di pubblicazione), lo storico, la gestione dei percorsi salvati/preferiti e la possibilità di eliminare singole tappe. Di seguito sono riportati i casi d'uso sviluppati per questa iterazione, in conformità con i requisiti iniziali.

\begin{table}[htbp]
\centering
\renewcommand{\arraystretch}{1.2}
\begin{tabular}{|l|l|}
\hline
\textbf{ID} & \textbf{Titolo} \\
\hline
UC3.1 & Visualizza catalogo \\
\hline
UC3.3 & Visualizza storico \\
\hline
UC4 & Filtri ricerca \\
\hline
UC5 & Salva percorso nei preferiti \\
\hline
UC7 & Rimuovi percorso salvato \\
\hline
UC8 & Pubblica percorso nel catalogo \\
\hline
UC10 & Imposta percorso completato \\
\hline
UC12.2 & Creazione itinerario casuale \\
\hline
UC15 & Modifica tappe \\
\hline
UC16 & Elimina tappa \\
\hline
\end{tabular}
\caption{Iterazione 2 - Casi d'uso implementati}
\label{tab:casi_uso_iterazione2}
\end{table}

\subsection*{UC3.1: Visualizza catalogo}
\textbf{Attori:} Utente, Sistema.\\
\textbf{Descrizione:} L'utente può visualizzare percorsi creati da altri utenti e resi pubblici, raccolti all'interno di un catalogo condiviso.\\
\textbf{Flusso degli eventi:}
\begin{enumerate}
    \item L'utente accede alla sezione "Catalogo" tramite l'interfaccia principale.
    \item Il sistema recupera dal database la lista dei viaggi con visibilità pubblica.
    \item Il sistema mostra all'utente l'elenco dei viaggi disponibili.
    \item L'utente può selezionare un singolo viaggio per visualizzarne i dettagli.
\end{enumerate}

\subsection*{UC3.3: Visualizza storico}
\textbf{Attori:} Utente, Sistema.\\
\textbf{Descrizione:} L'utente può consultare l'elenco dei propri viaggi passati, contrassegnati come completati.\\
\textbf{Flusso degli eventi:}
\begin{enumerate}
    \item L'utente accede alla sezione "Storico" dal proprio profilo.
    \item Il sistema interroga il database filtrando i viaggi dell'utente con stato \texttt{COMPLETED}.
    \item Il sistema restituisce e mostra a schermo l'elenco dei percorsi effettuati.
\end{enumerate}

\subsection*{UC4: Filtri ricerca}
\textbf{Attori:} Utente, Sistema.\\
\textbf{Descrizione:} L'utente può filtrare i percorsi nel catalogo in base a criteri quali città, durata massima o numero di tappe.\\
\textbf{Flusso degli eventi:}
\begin{enumerate}
    \item L'utente imposta i parametri del filtro nella sezione Catalogo.
    \item L'utente avvia la ricerca.
    \item Il sistema applica i filtri tramite una query su MongoDB.
    \item La vista viene aggiornata mostrando solo i risultati pertinenti.
\end{enumerate}

\subsection*{UC5: Salva percorso nei preferiti}
\textbf{Attori:} Utente, Sistema.\\
\textbf{Descrizione:} L'utente può aggiungere un percorso di interesse alla propria sezione dei preferiti.\\
\textbf{Flusso degli eventi:}
\begin{enumerate}
    \item L'utente visualizza un viaggio dal Catalogo o dai propri viaggi.
    \item Seleziona l'opzione di salvataggio nei preferiti.
    \item Il sistema crea l'associazione nel database.
    \item Il sistema conferma l'avvenuto salvataggio.
\end{enumerate}

\subsection*{UC7: Rimuovi percorso salvato}
\textbf{Attori:} Utente, Sistema.\\
\textbf{Descrizione:} L'utente può rimuovere un itinerario precedentemente salvato nella propria area personale.\\
\textbf{Flusso degli eventi:}
\begin{enumerate}
    \item L'utente accede alla sezione dei percorsi salvati.
    \item Seleziona il viaggio da rimuovere.
    \item Il sistema elimina l'associazione nel database.
    \item La lista viene aggiornata rimuovendo l'elemento dall'interfaccia.
\end{enumerate}

\subsection*{UC8: Pubblica percorso nel catalogo}
\textbf{Attori:} Utente, Sistema.\\
\textbf{Descrizione:} L'utente rende pubblico un proprio viaggio, permettendo ad altri utenti di visualizzarlo nel catalogo.\\
\textbf{Flusso degli eventi:}
\begin{enumerate}
    \item L'utente apre i dettagli di un proprio viaggio salvato.
    \item Seleziona l'opzione "Pubblica nel catalogo".
    \item Il sistema aggiorna lo stato di visibilità dell'itinerario nel database.
    \item Il percorso diventa visibile a tutti gli utenti nella sezione Catalogo.
\end{enumerate}

\subsection*{UC10: Imposta percorso completato}
\textbf{Attori:} Utente, Sistema.\\
\textbf{Descrizione:} L'utente contrassegna un viaggio come completato al termine della visita.\\
\textbf{Flusso degli eventi:}
\begin{enumerate}
    \item L'utente seleziona un viaggio attivo.
    \item Clicca su "Segna come completato".
    \item Il sistema aggiorna lo stato in \texttt{COMPLETED}.
    \item Il viaggio viene spostato automaticamente nella sezione Storico.
\end{enumerate}

\subsection*{UC12.2: Creazione itinerario casuale}
\textbf{Attori:} Utente, Sistema.\\
\textbf{Descrizione:} Generazione automatica di un itinerario basato su parametri casuali (città e numero tappe).\\
\textbf{Flusso degli eventi:}
\begin{enumerate}
    \item L'utente richiede un viaggio casuale specificando la città.
    \item Il sistema recupera i POI (da database o Overpass API).
    \item Il sistema seleziona un set casuale di monumenti rispettando i vincoli.
    \item Viene generato e mostrato il percorso ottimizzato sulla mappa.
\end{enumerate}

\subsection*{UC15: Modifica tappe}
\textbf{Attori:} Utente, Sistema.\\
\textbf{Descrizione:} L'utente modifica un viaggio esistente aggiungendo nuovi monumenti e richiedendo il ricalcolo dell'itinerario.\\
\textbf{Flusso degli eventi:}
\begin{enumerate}
    \item L'utente apre un viaggio in modalità modifica.
    \item Aggiunge una o più tappe selezionandole dalla mappa o dalla lista.
    \item Il sistema invoca l'Optimization Engine per ricalcolare il percorso TSP.
    \item Il viaggio aggiornato viene salvato con i nuovi dati di tempo e distanza.
\end{enumerate}

\subsection*{UC16: Elimina tappa}
\textbf{Attori:} Utente, Sistema.\\
\textbf{Descrizione:} L'utente può rimuovere una singola tappa da un itinerario esistente.\\
\textbf{Flusso degli eventi:}
\begin{enumerate}
    \item L'utente apre i dettagli di un viaggio salvato.
    \item Identifica la tappa da rimuovere e seleziona l'opzione di eliminazione.
    \item Il sistema rimuove la tappa dalla lista e ricalcola immediatamente l'ottimizzazione del percorso rimanente.
    \item L'interfaccia aggiorna la mappa e i dettagli del viaggio.
\end{enumerate}